% arXiv source. Compiles with pdflatex.
%
% GENERATED FILE, DO NOT EDIT. Source: paper/arxiv/main.tex.tmpl
% The placeholders are substituted from the measurement
% output by paper/arxiv/olcum/doldur.py, so that no number in the paper is
% typed by hand. Run paper/arxiv/olcum/olc.py first to produce the results
% file, then the substitution script to produce main.tex.
\pdfoutput=1

\documentclass[11pt]{article}

\usepackage[utf8]{inputenc}
\usepackage[T1]{fontenc}
\usepackage[margin=1in]{geometry}
\usepackage{amsmath}
\usepackage{amssymb}
\usepackage{amsthm}
\usepackage{booktabs}
\usepackage{graphicx}
\usepackage[hidelinks]{hyperref}

\newtheorem{proposition}{Proposition}
\newtheorem{remark}{Remark}

\title{Four Ways to Grow a Classifier and Why One of Them Cannot Learn}

\author{Cagri Temel\\
  \small Independent researcher\\
  \small \texttt{cagritemelusa@gmail.com}\\
  \small ORCID: 0009-0003-3359-6939}

\date{\today}

\begin{document}
\maketitle

\begin{abstract}
A classifier that builds itself has to decide where to put new capacity while it
trains. I measured four such decisions in tree structured and constructive
models: how to deepen a soft decision tree, how to initialise a new hidden unit,
which leaf to split next, and how to choose the model class at each node. Each is
usually presented as a free improvement. None of them is.

One cannot work at all, and I prove why. Deepen a soft decision tree so that its
function is unchanged, by turning every leaf into a gate whose two children
inherit the parent's class distribution, and the gradient with respect to the new
gate is identically zero. With the gate at one half the children receive
identical gradients, so they stay identical and the gradient stays zero. The
state reproduces itself at every step. This is a stronger degeneracy than the
symmetry Net2Net breaks with noise, or the saddle point that splitting steepest
descent escapes with second order information. Exact preservation costs
19.6 accuracy points on Iris, 19.1 on Wine and
55.6 on Digits against the same depth trained from scratch, and
inflates the fold to fold standard deviation. Any perturbation of the child
distributions removes the problem and its size barely matters.

The other three are trades. Fitting a new hidden unit to the residual before
installing it gives a smaller network on all three datasets but not a better one,
and on Digits it is significantly worse. Splitting the leaf with the most
expected error reached 0.885 with 3.7 splits where a
complete depth six tree needed 63, and lost 4.3 points on a harder
problem. Requiring statistical significance before a node gets a more expressive
split made the tree larger and less accurate.

Every number in this paper is written into it by the script that produced it.
\end{abstract}

\section{Introduction}

A constructive learner decides its own architecture while it trains. The idea is
old. Cascade correlation \cite{fahlman1990cascade} and the grow and learn
network \cite{alpaydin1994gal} both date from the early 1990s. It came back with
soft decision trees \cite{irsoy2012soft, frosst2017distilling}, which replace the
hard threshold at each internal node with a sigmoid gate and so make the whole
tree trainable by gradient descent. Once the tree is differentiable, growing it
during training is a design choice rather than a separate algorithm.

Most of those choices are settled by judgement. The candidate operations all
sound sensible: grow smoothly, keep what you have already learned, decide each
local question with a principled criterion. I implemented four of them, each
against the non growing baseline it is supposed to improve on, under one fixed
protocol.

The result I did not expect is how little the sensible sounding arguments
predict. Three of the four operations trade one quantity for another, and for
two of them the sign of the trade changes when the dataset changes. The fourth
does not trade anything, because it cannot learn at all. That one has a cause
that can be stated and proved rather than inferred from a table, and it is the
main contribution here.

I use \emph{grow} for any decision that adds structure while the model trains.
Three of the four operations add capacity: a level of the tree, a hidden unit, a
split at one leaf. The fourth chooses the model class at a node as the tree is
built. It is a construction decision rather than a capacity decision, and I
include it because it fails for a related reason.

The measurements came out of building \texttt{neural-trees}
\cite{temel2026neuraltrees}, a scikit-learn compatible \cite{pedregosa2011scikit}
library of these models with a PyTorch \cite{paszke2019pytorch} backend. I
should say plainly where they came from. Nothing here was designed in advance as
a four hypothesis study. These are the measurements that settled four
implementation questions, reported with the protocol held fixed and with the
negative outcomes left in. The negative ones have been more useful to me than
the positive ones, which is the reason for writing them up.

\section{Setup}

\subsection{Models}

A \emph{soft decision tree} of depth $d$ is a complete binary tree. Each
internal node $n$ carries a gate
\begin{equation}
  g_n(x) = \sigma(w_n^\top x + b_n), \qquad \sigma(z) = (1 + e^{-z})^{-1},
\end{equation}
which sends a fraction $g_n(x)$ of the probability mass arriving at $n$ to its
right child and $1 - g_n(x)$ to its left. Each leaf $\ell$ carries a class
distribution $Q_\ell \in \Delta^{K-1}$. Writing $\mu_\ell(x)$ for the product of
gate factors along the path from the root to $\ell$, the tree outputs
\begin{equation}
  \label{eq:mixture}
  p(y \mid x) = \sum_{\ell} \mu_\ell(x)\, Q_\ell ,
\end{equation}
and every parameter is trained by gradient descent on the cross entropy of
\eqref{eq:mixture}. Path products are accumulated in log space, so depth does not
cost precision.

A \emph{grow and learn (GAL) network} \cite{alpaydin1994gal} is a single hidden
layer network that adds a hidden unit when the training error stops improving and
removes units whose contribution has become negligible.

An \emph{omnivariate decision tree} \cite{yildiz2001omnivariate} chooses, at each
node independently, between a univariate threshold, a linear split and a
nonlinear split.

\subsection{Protocol}

Features are standardised to zero mean and unit variance using statistics from
the training fold only. Every number is computed over stratified five fold cross
validation repeated with three seeds, which is fifteen fits, and the same seeds
and folds are used for every arm of a comparison. I report the mean and, after
it, the standard deviation across those fifteen fits. That deviation describes
how much the quantity moves from fold to fold. It is not a confidence interval,
because the fifteen fits share training data and are not independent. Where a
difference needs a decision rather than a description I use the combined
$5\times2$cv $F$ test \cite{alpaydin1999combined}, which was designed for exactly
this situation.

Datasets are Iris ($n=150$, $p=4$, $K=3$), Wine ($n=178$, $p=13$, $K=3$), Breast
Cancer Wisconsin ($n=569$, $p=30$, $K=2$) and Digits ($n=1797$, $p=64$, $K=10$),
all as distributed with scikit-learn, plus two synthetic problems described where
they are used. Digits is included because the first three are small and
low dimensional, and a growth operation that only ever sees four features has not
been tested very hard.

No per arm hyperparameter search was performed. Both arms of each comparison use
identical settings apart from the operation under test, which isolates that
operation and deliberately gives up any claim that either arm is tuned to its
best.

\subsection{Reproducing this}

Every number in this paper is produced by one script,
\texttt{paper/arxiv/olcum/olc.py} in the repository, which writes a machine
readable results file. The manuscript is a template with named placeholders and
a second script substitutes the measured values into it. No number here was
typed by hand, which removes the failure mode where a table drifts away from the
code that produced it while nobody notices. The archived version is
\texttt{neural-trees} v0.6.2 \cite{temel2026neuraltrees}.

\section{Deepening a tree without disturbing it}

\subsection{The operation}

There is an obvious way to deepen a trained soft tree without losing what it has
learned. Turn every leaf $\ell$ into an internal node whose gate has all zero
parameters, so that $g(x) = \sigma(0) = 1/2$ everywhere, and give both new
children the parent's distribution, $Q_L = Q_R = Q_\ell$. The mass that used to
stop at $\ell$ now splits evenly between two identical children:
\begin{equation}
  \tfrac{1}{2} Q_L + \tfrac{1}{2} Q_R = Q_\ell ,
\end{equation}
so \eqref{eq:mixture} is unchanged. The deeper tree computes exactly the function
the shallower one computed, training continues from there, and the extra capacity
is available. This is the construction anyone would write down first, and it is
the tree analogue of the function preserving transformations of Net2Net
\cite{chen2016net2net} and network morphism \cite{wei2016network}.

\subsection{Why it cannot learn}

\begin{proposition}
\label{prop:zero}
Let a leaf of a soft decision tree be replaced by an internal node with gate
value $g$ and children carrying distributions $Q_L$ and $Q_R$, and let
$\mathcal{L}$ be any loss that depends on the parameters only through the tree
output \eqref{eq:mixture}.
\begin{enumerate}
\item If $Q_L = Q_R$ then $\partial \mathcal{L} / \partial w = 0$ and
      $\partial \mathcal{L} / \partial b = 0$ exactly, for any value of $g$.
\item If in addition $g = 1/2$ then
      $\partial \mathcal{L} / \partial Q_L = \partial \mathcal{L} / \partial Q_R$.
\end{enumerate}
Consequently a subtree initialised with $Q_L = Q_R$ and a zero gate stays in that
state under any update rule that is a deterministic function of the gradients
seen so far and is applied identically to identically situated parameters. The
added level never differentiates.
\end{proposition}

\begin{proof}
Write $o$ for the contribution of the subtree to the output,
$o = \mu_\ell \left( g\, Q_R + (1-g)\, Q_L \right)$, where $\mu_\ell$ is the
probability of arriving at the old leaf and does not depend on the new
parameters. Then
\begin{equation}
  \frac{\partial o}{\partial g} = \mu_\ell \left( Q_R - Q_L \right),
\end{equation}
which is the zero vector when $Q_L = Q_R$. The gate parameters influence
$\mathcal{L}$ only through $g$, so the chain rule gives
$\partial \mathcal{L}/\partial w = 0$ and $\partial \mathcal{L}/\partial b = 0$,
whatever $g$ happens to be. This proves the first claim. For the second,
\begin{equation}
  \frac{\partial \mathcal{L}}{\partial Q_L}
    = (1-g)\, \mu_\ell \frac{\partial \mathcal{L}}{\partial o},
  \qquad
  \frac{\partial \mathcal{L}}{\partial Q_R}
    = g\, \mu_\ell \frac{\partial \mathcal{L}}{\partial o},
\end{equation}
and the two coincide when $g = 1/2$. For the consequence, argue by induction on
the step count. At step $t$ the hypotheses hold, so the gate receives no update
and the two children receive the same update, which leaves $g = 1/2$ and
$Q_L = Q_R$ at step $t+1$. The hypotheses are restored exactly and the argument
repeats.
\end{proof}

\begin{remark}
The result does not depend on the loss, only on the output being the mixture
\eqref{eq:mixture}. It also does not depend on the gate being at one half: the
first claim holds for any $g$, and $g = 1/2$ is needed only to keep the children
moving together.
\end{remark}

\begin{remark}
Momentum, adaptive step sizes and weight decay do not rescue it. Every one of
them rescales or accumulates a gradient that is zero, and weight decay shrinks
two equal children by equal amounts. A batch does not rescue it either, because
both children see the same batch and therefore the same gradient.
\end{remark}

\begin{remark}
\label{rem:asimetri}
What does break it is any mechanism that makes the two children's updates differ.
Perturbing the child distributions at insertion is the direct route and the one
I measure below. Asymmetric noise during training is another: the subtree
dropout of \.Irsoy and Alpayd\i{}n \cite{irsoy2021dropout} drops branches
independently, which is exactly the condition the proposition rules out.

Floating point arithmetic does not supply such a perturbation here, and I
checked rather than assumed it. On a depth two tree trained on Iris and then
deepened, the gradient of every new gate is exactly zero in single precision,
the sum of absolute values over both weights and biases being $0$ and not
merely small, and the gradients of every pair of sibling leaves are bitwise
identical. Breaking the symmetry at insertion makes both quantities nonzero.
The check is in the repository and runs in continuous integration.
\end{remark}

\begin{remark}
The contrast with the neuron case is the point. Net2Net \cite{chen2016net2net}
duplicates a unit, notes that the copies are symmetric, and adds noise. Splitting
steepest descent \cite{liu2019splitting} shows that a duplicated neuron sits at a
saddle in the enlarged parameter space and computes the second order direction
that escapes it. GradMax \cite{evci2022gradmax} sets the new weights to zero but
gives the unit a bias, precisely so that the outgoing weights have a gradient to
follow. In all three the first order information is degenerate but recoverable.
In the tree gate case there is no direction to find at first order at all: the
derivative is not small, it is zero, and it reproduces itself. A second order
criterion in the style of \cite{liu2019splitting} would be the natural way to
repair the operation without noise, and I have not tried it.
\end{remark}

\subsection{What it costs}

Perturbing the child distributions at insertion time,
\begin{equation}
  Q_{L}, Q_{R} \leftarrow \mathrm{normalise}\!\left( Q_\ell \pm \varepsilon u \right),
  \qquad u \sim \mathcal{U}[0,1]^K ,
\end{equation}
breaks $Q_L = Q_R$ and with it both halves of Proposition~\ref{prop:zero}. The
function is then preserved only approximately. Table~\ref{tab:jitter} shows what
that approximation buys.

\begin{table}[ht]
\centering
\begin{tabular}{lccc}
\toprule
Growth & Iris & Wine & Digits \\
\midrule
Exactly function preserving ($\varepsilon = 0$) & 0.762 $\pm$ 0.103 & 0.786 $\pm$ 0.145 & 0.369 $\pm$ 0.035 \\
Perturbed, $\varepsilon = 0.05$ & 0.942 $\pm$ 0.041 & 0.979 $\pm$ 0.020 & 0.917 $\pm$ 0.034 \\
Perturbed, $\varepsilon = 0.2$ (the default) & 0.942 $\pm$ 0.044 & 0.979 $\pm$ 0.020 & 0.921 $\pm$ 0.027 \\
Perturbed, $\varepsilon = 0.5$ & 0.938 $\pm$ 0.050 & 0.983 $\pm$ 0.018 & 0.916 $\pm$ 0.026 \\
\midrule
Same depth, trained from scratch & 0.958 $\pm$ 0.032 & 0.978 $\pm$ 0.016 & 0.925 $\pm$ 0.029 \\
\bottomrule
\end{tabular}
\caption{Depth four soft decision tree grown one level at a time, as a function
of the perturbation applied to new child distributions. Mean and standard
deviation over three seeds of stratified five fold cross validation.}
\label{tab:jitter}
\end{table}

Three things in that table are worth separating.

The cost of exact preservation is large. It is 19.6 points on
Iris and 19.1 on Wine against the same depth trained from
scratch, and 55.6 on Digits.

The variance is worse than the mean suggests. At $\varepsilon = 0$ the fold to
fold standard deviation is 0.103 on Iris against
0.044 at the default perturbation. A model that is merely worse is
easier to live with than a model whose behaviour depends on which fold it landed
in, and the failure produces both.

The perturbation itself barely matters. Between $\varepsilon = 0.05$ and
$\varepsilon = 0.5$, a factor of ten, the means move by at most
0.4 points on any of the three datasets. Zero is not a small value
of $\varepsilon$. It is a different case, which is what
Proposition~\ref{prop:zero} says it should be.

Because this is the paper's central claim I test it rather than only describing
it. The combined $5\times2$cv $F$ test comparing $\varepsilon = 0$ against
$\varepsilon = 0.2$ gives $F = 4.33$, $p = 0.0596$ on Iris,
$F = 10.20$, $p = 0.0097$ on Wine and $F = 222.98$, $p < 0.0001$ on Digits.

The Iris result sits just above the conventional threshold and I am not going
to describe it as significant. It is the smallest dataset here, at 150 samples,
and the $5\times2$cv test is deliberately conservative. The other two reject
equality, and on Digits the margin is not close.

A regression test in the repository asserts that the gate gradient is zero when
the symmetry is not broken, so the mechanism itself is checked in continuous
integration rather than only in this paper.

\section{Where to put a new hidden unit}

The GAL network installs a new hidden unit with random incoming and outgoing
weights. Such a unit perturbs every output logit the moment it arrives, which
damages a network that has just been judged to have stopped improving, and it
then has to be trained from noise while the rest of the network repairs the
damage.

The alternative follows cascade correlation \cite{fahlman1990cascade}. Freeze
the network, compute its residual error on the training set, train the candidate
unit's incoming weights to correlate with that residual, and install the unit
with \emph{zero} outgoing weights. The network's function is unchanged at the
moment of insertion. That is safe here, unlike in Section~3, because the
outgoing weights have a nonzero gradient as soon as the incoming weights carry
any signal. GradMax \cite{evci2022gradmax} makes the same structural choice and
picks the incoming weights to maximise the gradient norm instead of to correlate
with the residual.

\begin{table}[ht]
\centering
\begin{tabular}{llccc}
\toprule
Dataset & New unit & Accuracy & Hidden units & $5\times2$cv $F$ \\
\midrule
Iris   & random   & 0.936 $\pm$ 0.046   & 17.2 $\pm$ 2.4   & \\
       & residual & 0.956 $\pm$ 0.037   & 6.9 $\pm$ 1.5   & \raisebox{1.4ex}[0pt]{$F = 2.33$, $p = 0.1811$} \\
\midrule
Wine   & random   & 0.985 $\pm$ 0.015   & 6.3 $\pm$ 2.2   & \\
       & residual & 0.985 $\pm$ 0.015   & 4.3 $\pm$ 0.5   & \raisebox{1.4ex}[0pt]{$F = 0.82$, $p = 0.6332$} \\
\midrule
Digits & random   & 0.949 $\pm$ 0.013 & 7.5 $\pm$ 0.6 & \\
       & residual & 0.928 $\pm$ 0.015 & 5.7 $\pm$ 0.5 & \raisebox{1.4ex}[0pt]{$F = 25.18$, $p = 0.0012$} \\
\bottomrule
\end{tabular}
\caption{GAL network, 150 epochs, three seeds of five fold cross validation.
The $F$ column tests the two accuracies against each other.}
\label{tab:gal}
\end{table}

The size effect is consistent and the accuracy effect is not, and I want to be
exact about which is which.

Every dataset gives a smaller network. Iris falls from 17.2 $\pm$ 2.4 hidden
units to 6.9 $\pm$ 1.5, Wine from 6.3 $\pm$ 2.2 to 4.3 $\pm$ 0.5,
Digits from 7.5 $\pm$ 0.6 to 5.7 $\pm$ 0.5. The mechanism is
visible in the growth trace. A random unit arrives unhelpful, the error does not
fall, and the growth criterion fires again, so the network grows because growth
is not working. Fitting the unit to the residual first removes that loop.

Accuracy goes the other way, and on the largest dataset it goes there
significantly. It rose on Iris and was identical on Wine to three decimals, and
the $F$ test rejects equality on neither. On Digits it fell by
2.1 points and the $F$ test does reject equality, in favour of
random initialisation.

So the operation is a parsimony intervention that costs accuracy on the one
dataset here with enough classes and features to strain it. The likely reason is
visible in the same table: the residual variant stops at 5.7 $\pm$ 0.5
hidden units where the random variant keeps going to 7.5 $\pm$ 0.6, and on
a ten class problem those extra units are doing work. Growth that stops early is
parsimony on an easy problem and underfitting on a harder one, and the growth
criterion cannot tell the two apart.

An earlier draft of this work reported the Iris accuracy gain on its own. That is
the kind of claim that survives exactly until somebody runs a fourth dataset.

One implementation detail cost me a day. The candidate objective involves a
normalisation of the form $s / \sqrt{v}$ where $v$ is a variance that can be zero
on a degenerate batch. At $v = 0$ the derivative of the square root is infinite
and the whole network becomes \texttt{NaN} in one step. Clamping $v$ away from
zero before the root fixes it. Clamping the result afterwards does not, because
by then the gradient has already been formed.

\section{Which leaf to split}

A soft tree can be grown a level at a time, as in Section~3, or a leaf at a time.
The per leaf rule of \.Irsoy, Y\i{}ld\i{}z and Alpayd\i{}n \cite{irsoy2012soft}
splits the leaf carrying the most expected error,
$\arg\max_\ell \sum_i \mu_\ell(x_i)\,e_i$, where $e_i$ is the error on sample
$i$.

\begin{table}[ht]
\centering
\begin{tabular}{llcc}
\toprule
Problem & Growth & Accuracy & Splits \\
\midrule
800 samples, 20 features  & complete depth 6 & 0.839 $\pm$ 0.021 & 63.0 $\pm$ 0.0 \\
                          & per leaf         & 0.885 $\pm$ 0.030   & 3.7 $\pm$ 1.4 \\
\midrule
2000 samples, 50 features & complete depth 6 & 0.579 $\pm$ 0.027 & 63.0 $\pm$ 0.0 \\
                          & per leaf         & 0.536 $\pm$ 0.040   & 4.8 $\pm$ 1.1 \\
\bottomrule
\end{tabular}
\caption{Growing one leaf at a time against a complete tree of the same maximum
depth, on two synthetic problems. Three seeds of five fold cross validation.}
\label{tab:perleaf}
\end{table}

On the first problem this is the result the rule is supposed to give. A complete
tree spends most of its capacity on regions that were already separated, and the
per leaf rule puts a split where the error is. It reached a higher accuracy with
roughly one seventeenth of the splits.

On the second problem the sparsity survives and the accuracy does not. The rule
still stops at 4.8 $\pm$ 1.1 splits, but it now loses 4.3 points to
the complete tree. Note also that both arms are poor there in absolute terms, so
what the second row shows is two ways of underfitting the same problem rather
than a good model beating a bad one.

I take the honest reading to be that the rule buys sparsity reliably and buys
accuracy only when the problem has enough structure for a handful of oblique
splits to capture it. Reported on the first problem alone it looks like a free
win. It is not one.

\section{Choosing the model class at a node}

The omnivariate tree picks each node's split type by cross validated accuracy on
that node's data. That rule will take a nonlinear split over a univariate one for
any improvement at all, however small and however likely to be noise. Requiring
the improvement to be statistically significant, by the combined $5\times2$cv $F$
test, and otherwise keeping the simpler split type, sounds strictly better:
simpler nodes at little or no cost.

\begin{table}[ht]
\centering
\begin{tabular}{llcc}
\toprule
Dataset & Selection rule & Accuracy & Nodes \\
\midrule
Breast Cancer & highest accuracy & 0.971 $\pm$ 0.018  & 3.8 $\pm$ 1.0 \\
              & significance     & 0.959 $\pm$ 0.020 & 8.1 $\pm$ 2.5 \\
\midrule
Digits        & highest accuracy & 0.510 $\pm$ 0.083  & 10.1 $\pm$ 1.7 \\
              & significance     & 0.492 $\pm$ 0.076 & 11.5 $\pm$ 2.3 \\
\bottomrule
\end{tabular}
\caption{Omnivariate tree of maximum depth three. Three seeds of five fold cross
validation.}
\label{tab:omni}
\end{table}

The rule does what it was asked to do at each node and the tree is worse for it
on both datasets. A simpler split separates its node's data less cleanly, so its
children inherit harder problems and have to be split in turn. Parsimony enforced
locally is paid for globally.

The absolute accuracies on Digits are low for a reason that has nothing to do
with the selection rule. A tree of maximum depth three has at most eight leaves
and Digits has ten classes, so the model cannot reach a high accuracy whatever it
does at each node. I kept the depth fixed across datasets rather than tuning it,
so the row is a valid comparison between the two rules and not a claim about how
well omnivariate trees classify digits.

There is a second problem, and it is structural rather than empirical. The test
needs a minimum sample count at a node to be meaningful, and nodes get smaller
with depth by construction. On the smaller datasets the rule almost never fires,
because almost every node below the root falls under the threshold. A
significance criterion is least available exactly where a recursive partitioning
method does most of its work. This is why the option ships switched off.

\section{What the three outcomes have in common}

The four outcomes line up on one axis: whether the optimiser has something to act
on at the moment the new structure arrives.

A unit fitted to the residual and installed with zero outgoing weights leaves the
function unchanged and has a nonzero gradient immediately. It reliably produces a
smaller network, which is what an operation that actually works looks like when
its benefit is parsimony. A split placed at the leaf with the most error changes
the function and puts capacity where the error already is. A level of gates
inserted with identical children leaves the function unchanged and has no
gradient at all, now or later, and it is the only one of the four that fails on
every dataset I tried.

The axis predicts which operations can work. It does not predict how much they
help, and Sections 4 and 5 show that the amount is dataset dependent enough to
change sign. Those are two different questions and I can only answer the first
one.

The practical rule is narrow and easy to check. When you add parameters to a
differentiable model during training, compute the gradient with respect to those
parameters on the first batch after insertion and assert that it is not zero. It
is one line in a test. In my case it was the difference between a growth
mechanism that worked and one that silently did nothing while consuming
parameters and compute.

\section{Limitations}

The datasets are small and standard, and Digits at $p=64$ is the largest thing
here. These are the problems the model families in question are normally
demonstrated on, but conclusions about accuracy should not be read as
generalising to large scale problems.

Each finding compares two arms of a single implementation with settings held
fixed. That isolates the operation under test. It does not establish that the
winning arm is the best available method, and it does not rule out an
implementation detail of ours being responsible for part of an effect.

Proposition~\ref{prop:zero} is exact and implementation independent. The
empirical results are neither. They are measurements of these operations in this
setting.

The standard deviations I report describe fold to fold spread over fifteen
non independent fits. They are not confidence intervals and should not be read as
such. Where I needed to decide rather than describe I used the combined
$5\times2$cv $F$ test, which is reported alongside.

I did not try the obvious repair suggested by Remark~\ref{rem:asimetri}, which
is to choose the split direction of a new gate from second order information in
the manner of splitting steepest descent \cite{liu2019splitting} instead of
perturbing the children at random. That is the experiment I would run next.

\section{Related work}

Soft decision trees are due to \.Irsoy, Y\i{}ld\i{}z and Alpayd\i{}n
\cite{irsoy2012soft}, and were revisited by Frosst and Hinton
\cite{frosst2017distilling} as a way of distilling a neural network into a model
whose decisions can be read. The hierarchical mixture of experts with subtree
dropout \cite{irsoy2021dropout} is the closely related architecture in which
leaves are themselves models. Omnivariate trees are due to Y\i{}ld\i{}z and
Alpayd\i{}n \cite{yildiz2001omnivariate}. Constructive learning of the kind
studied in Sections 4 and 5 goes back to cascade correlation
\cite{fahlman1990cascade} and GAL \cite{alpaydin1994gal}.

The idea of enlarging a network while preserving the function it computes is the
subject of Net2Net \cite{chen2016net2net} and network morphism
\cite{wei2016network}. Both note that an exactly preserving transformation
creates symmetric parameters, and Net2Net adds noise for that reason. Splitting
steepest descent \cite{liu2019splitting} treats the same situation as an
optimisation problem: a duplicated neuron sits at a saddle in the enlarged space,
and the paper derives the second order direction that leaves it, with a signed
refinement in later work. GradMax \cite{evci2022gradmax} keeps the function
unchanged by zeroing the new weights and then chooses the remaining freedom to
maximise the gradient norm, so that the new unit starts with the strongest
available learning signal.

Proposition~\ref{prop:zero} is the corresponding statement for tree gates, and it
is a stronger degeneracy than any of those. In the neuron case the first order
information is uninformative and second order information recovers a direction.
In the gate case, with identical children, the first derivative with respect to
the gate is exactly zero and the update rule reproduces the condition that made
it zero, so no amount of further first order training escapes it. The statistical
machinery in Section~6 is the combined $5\times2$cv $F$ test
\cite{alpaydin1999combined}.

\section{Conclusion}

Growth operations that look safest are not. Preserving a soft tree's function
exactly while deepening it produces a gate whose gradient is identically zero and
children that move together forever, and it costs around
19.6 points on Iris and 19.1 on Wine while
multiplying the fold to fold standard deviation by 1.3 to 7.3. Requiring statistical
significance before choosing a more expressive split makes each node simpler and
the tree larger and less accurate.

The two operations that worked share the property the failure lacks. Each adds
capacity that the optimiser can move on the next step. That is a cheap thing to
check, and it is now checked in a test.

\section*{Code and data availability}

The implementation, the measurement script that produced every number in this
paper, and the regression test asserting the zero gradient of
Proposition~\ref{prop:zero} are in the repository:

\begin{center}
\url{https://github.com/cgrtml/neural-trees}
\end{center}

\noindent The exact version used here is v0.6.2, archived at

\begin{center}
\url{https://doi.org/10.5281/zenodo.22802182}
\end{center}

\noindent All datasets are the standard versions distributed with scikit-learn.

\section*{Disclosure of AI assistance}

Parts of the software development and the drafting of this manuscript were
carried out with the assistance of a large language model used as a programming
and writing aid. Every empirical number reported here was produced by executing
the measurement script and substituting its output into this document
automatically. The author designed the experiments, verified the results and is
responsible for the content.

\begin{thebibliography}{99}

\bibitem{alpaydin1994gal}
E.~Alpayd\i{}n.
\newblock GAL: Networks that grow when they learn and shrink when they forget.
\newblock \emph{International Journal of Pattern Recognition and Artificial
  Intelligence}, 8(1):391--414, 1994.

\bibitem{alpaydin1999combined}
E.~Alpayd\i{}n.
\newblock Combined $5\times2$cv $F$ test for comparing supervised
  classification learning algorithms.
\newblock \emph{Neural Computation}, 11(8):1885--1892, 1999.

\bibitem{chen2016net2net}
T.~Chen, I.~Goodfellow, and J.~Shlens.
\newblock Net2Net: Accelerating learning via knowledge transfer.
\newblock In \emph{International Conference on Learning Representations}, 2016.

\bibitem{evci2022gradmax}
U.~Evci, B.~van Merri\"enboer, T.~Unterthiner, M.~Vladymyrov, and F.~Pedregosa.
\newblock GradMax: Growing neural networks using gradient information.
\newblock In \emph{International Conference on Learning Representations}, 2022.

\bibitem{fahlman1990cascade}
S.~E. Fahlman and C.~Lebiere.
\newblock The cascade-correlation learning architecture.
\newblock In \emph{Advances in Neural Information Processing Systems 2}, pages
  524--532, 1990.

\bibitem{frosst2017distilling}
N.~Frosst and G.~Hinton.
\newblock Distilling a neural network into a soft decision tree.
\newblock \emph{arXiv preprint arXiv:1711.09784}, 2017.

\bibitem{irsoy2012soft}
O.~\.Irsoy, O.~T. Y\i{}ld\i{}z, and E.~Alpayd\i{}n.
\newblock Soft decision trees.
\newblock In \emph{Proceedings of the 21st International Conference on Pattern
  Recognition}, pages 1819--1822, 2012.

\bibitem{irsoy2021dropout}
O.~\.Irsoy and E.~Alpayd\i{}n.
\newblock Dropout regularization in hierarchical mixture of experts.
\newblock \emph{Neurocomputing}, 419:148--156, 2021.

\bibitem{liu2019splitting}
Q.~Liu, L.~Wu, and D.~Wang.
\newblock Splitting steepest descent for growing neural architectures.
\newblock In \emph{Advances in Neural Information Processing Systems 32}, 2019.

\bibitem{paszke2019pytorch}
A.~Paszke et~al.
\newblock PyTorch: An imperative style, high-performance deep learning library.
\newblock In \emph{Advances in Neural Information Processing Systems 32}, pages
  8024--8035, 2019.

\bibitem{pedregosa2011scikit}
F.~Pedregosa et~al.
\newblock Scikit-learn: Machine learning in Python.
\newblock \emph{Journal of Machine Learning Research}, 12:2825--2830, 2011.

\bibitem{temel2026neuraltrees}
C.~Temel.
\newblock neural-trees: soft decision trees, hierarchical mixtures of experts
  and classifier comparison tests for Python, version 0.6.2.
\newblock Zenodo, 2026. \url{https://doi.org/10.5281/zenodo.22802182}

\bibitem{wei2016network}
T.~Wei, C.~Wang, Y.~Rui, and C.~W. Chen.
\newblock Network morphism.
\newblock In \emph{Proceedings of the 33rd International Conference on Machine
  Learning}, pages 564--572, 2016.

\bibitem{yildiz2001omnivariate}
O.~T. Y\i{}ld\i{}z and E.~Alpayd\i{}n.
\newblock Omnivariate decision trees.
\newblock \emph{IEEE Transactions on Neural Networks}, 12(6):1539--1546, 2001.

\end{thebibliography}

\end{document}
